\begin{theorem}[Smith \texorpdfstring{$p$}{p}-主谱的 Ramanujan 素数幂影公式]\label{thm:conclusion-smith-pprimary-ramanujan-shadow-formula}
设 \(A\in M_{m\times n}(\ZZ)\) 满足 \(\operatorname{rank}_{\QQ}(A)=m\)，其 Smith 不变量记为
\[
d_1\mid d_2\mid\cdots\mid d_m.
\]
对素数 \(p\)，记
\[
e_i(p):=v_p(d_i),
\qquad
\Delta_p(k):=\#\{\,i:e_i(p)\ge k\,\}
\quad (k\ge 1),
\qquad
\Delta_p(0):=m,
\]
并定义 prime-power Ramanujan 影
\[
\mathcal R_{A,p}(k):=\sum_{i=1}^{m}c_{p^k}(d_i)
\qquad (k\ge 1).
\]
则对任意 \(k\ge 1\)，成立精确恒等式
\[
\mathcal R_{A,p}(k)=p^k\Delta_p(k)-p^{k-1}\Delta_p(k-1).
\]
因此，单个素数塔上的 Ramanujan 影精确记录了 \(p\)-主 Smith 图样的列高差分。
\leanverified{shadow\_formula\_algebraic}
\leanverified{shadow\_formula\_telescoping\_K3}
\leanverified{paper\_conclusion\_smith\_ramanujan\_shadow\_formula}
\end{theorem}
\begin{proof}
对素数幂导子，Ramanujan 和满足
\[
c_{p^k}(n)=
\begin{cases}
0, & v_p(n)\le k-2,\\[2pt]
-p^{k-1}, & v_p(n)=k-1,\\[2pt]
p^{k-1}(p-1), & v_p(n)\ge k.
\end{cases}
\]
因此
\[
\mathcal R_{A,p}(k)
=
p^{k-1}(p-1)\Delta_p(k)-p^{k-1}\bigl(\Delta_p(k-1)-\Delta_p(k)\bigr),
\]
整理即得
\[
\mathcal R_{A,p}(k)=p^k\Delta_p(k)-p^{k-1}\Delta_p(k-1).
\]
证毕。
\end{proof}

\begin{theorem}[prime-power Ramanujan 影的完全反演与实现判别]\label{thm:conclusion-smith-pprimary-ramanujan-shadow-inversion}
固定素数 \(p\) 与整数 \(m\ge 1\)。对任意数列 \((R_k)_{k\ge 1}\)，定义
\[
\Delta_k:=p^{-k}\left(m+\sum_{j=1}^{k}R_j\right)
\qquad (k\ge 1).
\]
则 \((R_k)_{k\ge 1}\) 可表示成某个满行秩整数矩阵的 \(p\)-主 Ramanujan 影
\[
R_k=\mathcal R_{A,p}(k)
\qquad (k\ge 1)
\]
当且仅当 \((\Delta_k)_{k\ge 1}\) 是一个非负整数、单调不增并最终为零的序列。并且在该情形下，
\[
\#\{\,i:e_i(p)=k\,\}=\Delta_k-\Delta_{k+1}.
\]
\leanverified{paper_conclusion_smith_ramanujan_shadow_inversion_seeds}
\leanverified{paper_conclusion_smith_ramanujan_shadow_inversion_package}
\leanverified{paper\_conclusion\_smith\_pprimary\_ramanujan\_shadow\_inversion}
\end{theorem}
\begin{proof}
若 \(R_k=\mathcal R_{A,p}(k)\)，则由定理 \ref{thm:conclusion-smith-pprimary-ramanujan-shadow-formula}，
\[
\mathcal R_{A,p}(k)=p^k\Delta_p(k)-p^{k-1}\Delta_p(k-1).
\]
令 \(E_k:=p^k\Delta_p(k)\)，则
\[
\mathcal R_{A,p}(k)=E_k-E_{k-1},
\qquad
E_0=\Delta_p(0)=m.
\]
望远镜求和即得
\[
E_k=m+\sum_{j=1}^{k}\mathcal R_{A,p}(j),
\]
亦即
\[
\Delta_p(k)=p^{-k}\left(m+\sum_{j=1}^{k}\mathcal R_{A,p}(j)\right).
\]
因此 \((\Delta_k)\) 必为非负整数、单调不增且最终为零。再由列高差分公式，
\[
\#\{\,i:e_i(p)=k\,\}=\Delta_p(k)-\Delta_p(k+1).
\]

反之，若 \((\Delta_k)\) 满足上述三个条件，则
\[
m_k:=\Delta_k-\Delta_{k+1}\in \ZZ_{\ge 0},
\qquad
\sum_{k\ge 0}m_k=m.
\]
取恰有 \(m_k\) 个 \(p\)-主赋值等于 \(k\) 的多重集 \(\{e_i(p)\}_{i=1}^{m}\)，再取对应 Smith 对角矩阵即可实现该多重集。按定理 \ref{thm:conclusion-smith-pprimary-ramanujan-shadow-formula} 反算，即得 \(R_k=\mathcal R_{A,p}(k)\)。证毕。
\end{proof}

\begin{theorem}[局部 zeta 因子与 Ramanujan 影的双向有限差分同构]\label{thm:conclusion-smith-local-zeta-ramanujan-shadow-difference-equivalence}
在上述记号下，定义
\[
f_p(-1):=-m,
\qquad
f_p(0):=0,
\qquad
f_p(k):=\sum_{i=1}^{m}\min(k,e_i(p))
\quad (k\ge 1).
\]
则对一切 \(k\ge 1\)，成立二阶有限差分恒等式
\[
\mathcal R_{A,p}(k)
=
p^{k-1}\Bigl(
p\,f_p(k)-(p+1)f_p(k-1)+f_p(k-2)
\Bigr).
\]
反过来，
\[
f_p(k)
=
\sum_{j=1}^{k}\Delta_p(j)
=
\sum_{j=1}^{k} p^{-j}\left(m+\sum_{\ell=1}^{j}\mathcal R_{A,p}(\ell)\right).
\]
因而，局部因子
\[
Z_{A,p}(s)=\sum_{k\ge 0}p^{-k(s-n+m)+f_p(k)}
\]
与 Ramanujan 影序列 \((\mathcal R_{A,p}(k))_{k\ge 1}\) 彼此唯一决定。
\leanverified{paper\_conclusion\_smith\_local\_zeta\_ramanujan\_shadow\_difference\_equivalence}
\end{theorem}
\begin{proof}
由定义，
\[
f_p(k)-f_p(k-1)=\Delta_p(k)
\qquad (k\ge 1).
\]
把它代入定理 \ref{thm:conclusion-smith-pprimary-ramanujan-shadow-formula} 的
\[
\mathcal R_{A,p}(k)=p^k\Delta_p(k)-p^{k-1}\Delta_p(k-1),
\]
便得
\[
\mathcal R_{A,p}(k)
=
p^{k-1}\Bigl(
p(f_p(k)-f_p(k-1))-(f_p(k-1)-f_p(k-2))
\Bigr),
\]
即所述二阶差分公式。

另一方面，由上一条定理，
\[
\Delta_p(j)=p^{-j}\left(m+\sum_{\ell=1}^{j}\mathcal R_{A,p}(\ell)\right),
\]
对 \(j=1,\dots,k\) 求和即得 \(f_p(k)\) 的反演公式。最后把 \((f_p(k))\) 代入局部 zeta 因子闭式，便知 \((\mathcal R_{A,p}(k))\) 与 \(Z_{A,p}\) 彼此唯一决定。证毕。
\end{proof}

\begin{proposition}[Smith 光谱的 Ramanujan--Dirichlet 阴影]\label{prop:conclusion-smith-global-ramanujan-dirichlet-shadow}
\leanverified{paper\_conclusion\_smith\_global\_ramanujan\_dirichlet\_shadow}
定义全局 Ramanujan 影
\[
\mathcal R_A(q):=\sum_{i=1}^{m}c_q(d_i),
\]
以及其 Dirichlet 级数
\[
\mathfrak R_A(s):=\sum_{q\ge 1}\frac{\mathcal R_A(q)}{q^s}.
\]
则当 \(\Re s>1\) 时，
\[
\mathfrak R_A(s)=\frac{1}{\zeta(s)}\sum_{i=1}^{m}\sigma_{1-s}(d_i).
\]
\leanverified{paper\_conclusion\_smith\_global\_ramanujan\_dirichlet\_shadow}
\end{proposition}
\begin{proof}
经典恒等式给出
\[
\sum_{q\ge 1}\frac{c_q(n)}{q^s}=\frac{\sigma_{1-s}(n)}{\zeta(s)}
\qquad (\Re s>1).
\]
对 \(n=d_i\) 求和并交换有限求和次序，即得
\[
\mathfrak R_A(s)
=
\sum_{i=1}^{m}\sum_{q\ge 1}\frac{c_q(d_i)}{q^s}
=
\frac{1}{\zeta(s)}\sum_{i=1}^{m}\sigma_{1-s}(d_i).
\]
证毕。
\end{proof}

\begin{theorem}[有限 Ramanujan 审计的不完备性]\label{thm:conclusion-smith-truncated-ramanujan-audit-incompleteness}
固定有限素数集 \(S\subset \PP\) 与截断高度 \(K\ge 1\)。定义截断 Ramanujan 审计数据
\[
\mathfrak R_{S,K}(A):=\{\mathcal R_{A,p}(k):p\in S,\ 1\le k\le K\}.
\]
则 \(\mathfrak R_{S,K}\) 不是完整 Smith 光谱的完备不变量；更强地说，每一个固定的 \(\mathfrak R_{S,K}\)-等价类都包含无穷多个两两不同的 Smith 光谱。
\leanverified{paper\_conclusion\_smith\_truncated\_ramanujan\_audit\_incompleteness}
\end{theorem}
\begin{proof}
由定理 \ref{thm:conclusion-truncated-smith-ledger-infinite-indistinguishability}，对任意固定的 \(S\) 与 \(K\)，都存在无穷多个满行秩整数矩阵 \(A^{(r)}\) 使得
\[
K_p(k;A^{(r)})=K_p(k;A)
\qquad
(p\in S,\ 1\le k\le K),
\]
而它们的完整 Smith 不变量系两两不同。
由局部核公式，
\[
K_p(k;A)=p^{k(n-m)+f_{A,p}(k)},
\]
故这些矩阵在同一有限窗口内具有完全相同的 \(f_{A,p}(k)\)。再由定理 \ref{thm:conclusion-smith-local-zeta-ramanujan-shadow-difference-equivalence}，它们在该窗口内的 \(\mathcal R_{A,p}(k)\) 亦完全相同。于是 \(\mathfrak R_{S,K}\) 不能唯一决定完整 Smith 光谱。证毕。
\end{proof}

\begin{theorem}[prime-power Ramanujan 审计的精确 cutoff 完备性]\label{thm:conclusion-smith-pprimary-ramanujan-cutoff-completeness}
\leanverified{paper\_conclusion\_smith\_pprimary\_ramanujan\_cutoff\_completeness}
对固定素数 \(p\)，记
\[
E_p^{\max}:=\max_{1\le i\le m}e_i(p).
\]
则截断序列
\[
\bigl(\mathcal R_{A,p}(1),\dots,\mathcal R_{A,p}(K)\bigr)
\]
唯一决定 \(A\) 的 \(p\)-主 Smith 光谱，当且仅当
\[
K\ge E_p^{\max}.
\]
\end{theorem}
\begin{proof}
若 \(K\ge E_p^{\max}\)，则对所有 \(k>K\) 有 \(\Delta_p(k)=0\)。由定理 \ref{thm:conclusion-smith-pprimary-ramanujan-shadow-inversion}，
\[
\Delta_p(k)=p^{-k}\left(m+\sum_{j=1}^{k}\mathcal R_{A,p}(j)\right)
\]
已由截断数据完全确定；而当 \(k>K\) 时自动为零。因此全部 \(\Delta_p(k)\) 都可恢复，进而
\[
\#\{\,i:e_i(p)=k\,\}=\Delta_p(k)-\Delta_p(k+1)
\]
给出完整 \(p\)-主光谱。

反之，若 \(K<E_p^{\max}\)，则存在至少一个 \(i\) 使 \(e_i(p)>K\)。把该指数从 \(E\) 改成任意更大的 \(E'>E\)，则对所有 \(1\le k\le K\)，\(\Delta_p(k)\) 保持不变，从而按定理 \ref{thm:conclusion-smith-pprimary-ramanujan-shadow-formula}，
\[
\mathcal R_{A,p}(k)
\qquad
(1\le k\le K)
\]
亦保持不变；然而 \(p\)-主 Smith 光谱已经改变。故当 \(K<E_p^{\max}\) 时，截断序列不可能完备。证毕。
\end{proof}

\begin{conjecture}[Ramanujan prime-power 影的初对象性]\label{conj:conclusion-smith-ramanujan-shadow-initiality}
在一切满足以下三条的 \(p\)-局部不变量族中：
\begin{enumerate}
\normalfont
  \item 仅依赖于 \(\{e_i(p)\}_{i=1}^{m}\)；
  \item 对 block sum 可加，或等价地可经有限差分递推恢复；
  \item 对有限 truncation 具有精确 cutoff 完备门槛，
\end{enumerate}
prime-power Ramanujan 影族
\[
\{\mathcal R_{A,p}(k)\}_{k\ge 1}
\]
应当是初对象：任何其他此类完备族，都应经由一个上三角、有限差分型的有理变换从 \(\{\mathcal R_{A,p}(k)\}\) 导出。

其根据是刚性的：定理 \ref{thm:conclusion-smith-pprimary-ramanujan-shadow-formula} 表明它与列高函数 \(\Delta_p(k)\) 一阶望远镜等价，定理 \ref{thm:conclusion-smith-local-zeta-ramanujan-shadow-difference-equivalence} 表明它与局部 zeta 因子二阶有限差分等价，而定理 \ref{thm:conclusion-smith-pprimary-ramanujan-cutoff-completeness} 又给出恰好由 \(E_p^{\max}\) 控制的精确 cutoff 门槛。因此，若还存在另一族同等完备的 \(p\)-局部坐标，它应当只是 Ramanujan 影的结构性重写，而不是新的原子层。
\end{conjecture}
